\section{Conclusion Draft}
\label{sec:conclusion}

This synthesis organizes completed Double Heston inverse-calibration evidence around a
single defensible message: a validated forward engine and a frozen known-truth dataset
make it possible to compare fast learned inverses against a strong traditional frontier
without conflating speed, structural validity, repricing fit, and parameter recovery.

Under that discipline, traditional numerical calibration remains the accuracy reference
but is computationally expensive. Ordinary ANN and constraint/repricing-informed neural
mappings are dramatically faster and structurally valid, with moderate aggregate errors
and excellent seed stability, but materially larger known-truth errors. Their similar
performance indicates that the differentiable repricing term did not solve the central
inverse problem in the frozen configuration.

The most important conclusion is negative and constructive: low repricing error can
coexist with material parameter displacement. Any useful calibration methodology should
therefore report tolerance-conditioned parameter recovery separately from fit-to-
observations, retain failed starts, disclose constraints and starts, and avoid claiming
unique identification from optimizer success.

The next evidentiary steps are already bounded: finish positive-noise traditional subset
calibration under its frozen protocol, evaluate genuine PDE-informed Model3 only after
the frozen Stage A/B sequence, and reserve OOD and final real-market conclusions for their
separately authorized experiments. Until then, the canonical claim is a frozen
representation plus retained practical non-identifiability, not a solved inverse problem.
